\section{Discussion}
\label{sec:discussion}

\subsection{What geometric calibration constrains}
A length calibration constrains the combination of source strength, constitutive assumptions and observation operator that reproduces one selected boundary extent. In the present model, agreement with B length leaves a wider and shallower fusion envelope than the reference, and the relative contraction of width and depth differs from the nominal B-to-C response (Fig.~\ref{fig:comparison}). The discrepancy therefore concerns energy distribution as well as its overall scale. Fitting the source factor to length does not supply an independent constraint on penetration, rear phase separation or the underlying temperature profile.

This distinction has a direct precedent in Myers et al.~\cite{myers2023}. Their bare-plate 316L FLOW-3D calculations fitted Fresnel and accommodation coefficients simultaneously to cross-sectional area and width-to-depth ratio across six power--speed combinations. Parameter combinations with similar geometric agreement produced different surface temperatures; two-color profiles provided an additional discriminating observable. That model included melt motion and vaporization but used constant surface tension. Its finding establishes a geometry--temperature ambiguity under those conditions, rather than attributing the present residual to the same parameters. Here only the baseline length is fitted, and the continuation branches retain its source factor. Their differences consequently test response after one calibration, rather than demonstrate different thermal histories among equally refitted models.

Kollmannsberger et al.~\cite{kollmannsberger2019} provide the closer AMMT geometry comparison. With the measured laser input constraining the source shape, scalar absorption calibration left appreciable width/depth errors; effective directional conductivity calibrated at B improved the A/C geometry comparisons. Their temperature-dependent, latent-inclusive formulation already addresses nonlinear material behavior. The contrast with the present isotropic Gaussian calculation shows that source representation and directional closure can change geometric agreement. It does not rank the solvers, because the beam, material continuation, phase treatment and extraction differ. Importantly, their study explicitly limited the validity of internal melt-pool temperatures despite improved geometry.

Recent reliability studies reinforce the need to distinguish observables. Soares et al.~\cite{soares2026} compared a Gaussian conduction model with thermal-camera and metallographic measurements for powder tracks in 316L, IN625 and CoCr. Properties were fixed near $0.8T_{\rm melt}$, with powder corrections and no explicit phase change. Good correlation coexisted with biased surface area and regime-dependent depth discrepancy. Their camera calibration used melting and metallographic information, so these channels were not wholly independent. The present bare-plate enthalpy model asks a different question, but shares the lesson that agreement in one dimension or correlation across conditions does not establish agreement of the full field.

\subsection{The information gained by fixing the calibrated source}
Keeping the source fixed makes the closure contrast interpretable: each change is a response to a specified continuation choice, without an absorption adjustment absorbing part of it. The joint length increase of 34.28~\um{} combines a conductivity contribution of +43.21~\um{}, a sensible-storage contribution of $-6.17$~\um{} and an interaction of $-2.76$~\um{}. Conductivity holding lengthens the pool at both storage settings, whereas storage holding shortens it at both conductivity settings. A joint property replacement would conceal these opposing responses. The small joint depth change similarly conceals cancellation, and none of the four corners repairs the B wide/shallow pattern (Table~\ref{tab:factorial}).

The rear profiles explain what the length response represents. Both rear phase boundaries move together by approximately 43~\um{} under conductivity holding, while their separation changes by approximately 1~\um{} (Fig.~\ref{fig:factorial}). Rear solidus displacement supplies 98.3\% of the length change. Thus the dominant response is a translation of the thermal tail relative to the source, rather than a comparable enlargement of the axial phase interval. Tracking solidus and liquidus separately exposes information that total length alone suppresses.

Material sensitivity itself is established. Yang et al.~\cite{yang2024} varied constant properties and source parameters in analytical Ti6Al4V models and assessed their contributions to geometry and temperature. Whalen et al.~\cite{whalen2021} incorporated temperature-dependent properties and powder conditions into an adapted Eagar--Tsai model, inferred distributions of apparent absorption and porosity, and propagated geometry uncertainty for 316L. The present deterministic four-corner comparison instead separates two missing-range functions in one nonlinear enthalpy model and connects their finite interaction to rear-threshold motion. Its effect magnitudes are conditional responses, rather than experimental property estimates or probabilistic uncertainty bounds.

\subsection{Transport and storage: a conditional mechanism}
Latent-inclusive storage helps explain the modest sensitivity to sensible heat capacity in the adopted phase interval. The fixed linear phase law adds 4666.67~\si{\joule\per\kilogram\per\kelvin} to effective specific heat between solidus and liquidus. At 1320~$^\circ$C, holding sensible $c_p$ reduces it by 7.09\%, but reduces latent-inclusive effective specific heat by only 0.95\%; enthalpy accumulated from the cold reference changes by 0.66\%. Conductivity decreases by 22.48\% at the same temperature. These are distinct local comparisons: the sensible-property contrast is substantially moderated when expressed as total storage through the phase interval.

This supports a storage-based interpretation of the unequal responses without making latent heat an independently isolated cause. Latent heat and phase endpoints are unchanged in every corner, and the $c_p$ continuation begins below the solidus. It therefore also changes sensible enthalpy acquired before melting and the temperature--enthalpy inverse throughout the hotter field. The observed shortward displacement under storage holding is consistent with a changed enthalpy-transport balance, but a matched-temperature percentage alone cannot predict the global sign or magnitude. The relative sensitivities could differ with another freezing interval, phase law, source or property function; nonlinear enthalpy treatments such as Van Elsen et al.~\cite{vanelsen2007} likewise make phase effects conditional on the transport setting.

Conductivity holding reduces local diffusion relative to coordinate transport in the discretized mushy region, consistent with the retained axial face P\'eclet median increasing from 4.44 to 5.77. This transport is material passage through the laser frame, not liquid convection. Moreover, the selected face population changes with the field. The diagnostic supports a shift in the model's transport balance, rather than isolating an exclusive bulk mechanism.

The integrated flux supplies a useful counterexample to a simpler explanation. Outward diffusive power from the selected rear hot region increases from 14.724 to 14.913~W under conductivity holding. Lower conductivity therefore does not imply lower integrated outward heat flow in this comparison: temperature gradients and the temperature-selected region boundary change simultaneously. Conductivity also enters surface recovery through its primitive and inverse, so the reported surface crossings contain both the solved field response and a consistently changed recovery step. The large tail displacement survives as an observation of this model; interpreting it solely as suppressed bulk heat escape would exceed the available diagnostics.

\subsection{Spatial phase structure and material time}
The nearly equal B/C rear phase spans, 80.26 and 80.79~\um{}, coexist with passage times of 100.33 and 67.33~\us{}. The distinction follows from dividing spatial separation by scan speed: faster passage shortens the material interval even when the modeled spatial phase structure changes little. The associated 49.02\% increase in interval-mean cooling rate is therefore a consequence of the same thresholds and coordinate transformation, rather than another independent validation endpoint. It also differs from length sensitivity: moving both rear thresholds together changes their location relative to the laser without necessarily changing their separation.

This distinction matters when comparing simulated fields with thermography. Hooper~\cite{hooper2018} measured spatial profiles and fixed-location thermal histories in Ti6Al4V, including scan events with significant transients. Lane et al.~\cite{lane2020} provided cooling-rate exemplars using a below-solidus interval and an imaging-based rear-boundary definition. Neither is a matched measurement of the present 1350--1290~$^\circ$C centerline interval. A physical comparison would require a temperature trajectory over the same interval, with scan position, exposure and observation mapping accounted for. The transformed interval is informative about this quasi-steady model but does not determine segregation, grain morphology or the alloy's actual nonequilibrium freezing path.

\subsection{Surviving explanations and discriminating comparisons}
Three explanations remain compatible with the shape discrepancy. First, the Gaussian surface source and constant effective factor may distribute energy incorrectly across the conditions. Calibrated conduction models can transfer successfully: Ross et al.~\cite{ross2022} inferred volumetric source parameters from solidification-boundary temperatures and interpolated their trends; Kusano and Watanabe~\cite{heatsourcebayes2024} compared source families and constructed a calibrated conical-source regression for an IN738LC process window. Coleman et al.~\cite{dynamic2024} calibrated a dynamic volumetric source for bare IN625 at 195~W and 800~mm/s over multiple spot sizes, then applied it to an AMB2018-01 layer case. Their phase-dependent properties and 1410~K eutectic endpoint differ from the present equilibrium interval. These successes make source choice a substantive alternative to changing conductivity alone.

Second, constitutive uncertainty and directional melt transport may affect shape and rear temperature differently. The tested scalar continuations strongly influence length while leaving the main B shape residual. This restricts those two scenarios as geometric repairs, but does not eliminate other high-temperature functions or effective directional transport. Composition-based estimates can provide physically motivated candidate functions where measurements are unavailable~\cite{mills2006}; they cannot identify this coupon's liquid properties. A comparison using the same measured beam, phase assumptions, length target and extraction operators would clarify whether an alternative source or directional closure improves excluded width/depth together with the rear profile. Adding an independently calibrated thermal profile would distinguish parameter compensation from credible thermal agreement.

Third, observation mapping or omitted free-surface physics may contribute. The experimental radiance length and the modeled solidus-crossing proxy differ, while fusion width/depth use different populations and a passage-maximum operator. An imaging replay could assess how this mismatch affects fitted source strength before introducing additional physics. Melt-flow studies show how recoil, surface-tension-driven motion and evaporative cooling redistribute energy~\cite{khairallah2016}, and X-ray imaging directly reveals vapor-depression transitions~\cite{cunningham2019}. Those results establish viable mechanisms under their own conditions, rather than diagnose a keyhole or flow mechanism from the present wide/shallow residual.

Hong et al.~\cite{hong2026} illustrate a more discriminating failure test. Their Eagar--Tsai bare-track analysis uses melting-range effective properties, separately specified absorption inputs, and depth inversion to ask whether any physical absorption can reach the measured penetration; property and directional-diffusivity controls test that criterion. IN625 absorption uses an empirical relation rather than a matched in-situ measurement. The present study performs neither that inversion nor those source/transport controls. Its shape discrepancy therefore identifies a model limitation after fitting, without establishing the physical necessity of a vapor cavity. The useful next comparison is one that constrains source/measurement mapping and tests an excluded thermal or shape observable, rather than selecting a mechanism from the residual alone.

\subsection{Limits on thermal interpretation}
Two constraints govern the strength of this diagnosis. The continuation functions start above supplier-table endpoints of 982~$^\circ$C for $k$ and 1093~$^\circ$C for $c_p$, both below the solidus~\cite{specialmetals625}. Neither is measured liquid behavior, and the equilibrium phase law remains assumed. Held-conductivity peak temperatures are sensitivity outputs outside credible physical-temperature prediction. Numerical verification, including the simpler analytical benchmark and retained axial/domain comparisons, supports the scale of the main length and shape contrasts~\cite{oberkampf2002}; it supplies neither a full three-dimensional error bound for every branch nor precision for submicrometre interactions.

The experimental comparisons are retrospective and nonblind, with one fitted length, excluded shape diagnostics and separately sourced uncertainty scales. They cannot form a joint statistical validation of the thermal field. The earned implication is consequently specific: geometric fitting remains useful, but closure response, boundary translation, phase separation and material time must be distinguished before assigning solidification meaning to the fitted solution.
